\documentclass[12pt,a4paper]{article}

\usepackage[utf8]{inputenc}
\usepackage[T1]{fontenc}
\usepackage{amsmath,amssymb,amsthm}
\usepackage{doi}
\usepackage{hyperref}

\newtheorem{theorem}{Theorem}[section]
\newtheorem{lemma}[theorem]{Lemma}
\newtheorem{corollary}[theorem]{Corollary}
\theoremstyle{remark}
\newtheorem{remark}[theorem]{Remark}
\theoremstyle{definition}
\newtheorem{definition}[theorem]{Definition}

\newcommand{\SO}{\mathrm{SO}}
\newcommand{\Meff}{M_{\mathrm{eff}}}

\title{Lorentz Transformations from Projective Temporal Ordering\\[4pt]
\normalsize\itshape LC-O1 of the Lorentz--Capacity Sub-Programme}
\author{J\'er\^ome Beau\\[2pt]
\small Independent Researcher\\
\small \texttt{jerome.beau@cosmochrony.org}}
\date{2026}

\begin{document}

  \maketitle

  \begin{abstract}
    We close LC-O1, the reconstruction of effective Lorentz transformations between
    inertial observers from projective temporal ordering data.
    The Lorentz-capacity construction of~\cite{Beau2026n,Beau2026tp} establishes that
    the projective temporal capacity satisfies
    \[
      F^\tau = \sqrt{1-\beta^2} = \frac{1}{\gamma},
    \]
    where $\beta$ is the normalised spatial capacity load and $\gamma$ is the Lorentz
    factor, closing time dilation and the light cone.
    The remaining open problem is the reconstruction of the full observer transformation,
    including length contraction.
    We prove that an inertial observer corresponds to a unit timelike splitting of the
    completed effective co-metric $g^{\mu\nu}=2\eta^{\mu\nu}$ established in
    Q5b--Q11~\cite{Beau2026q5b,Beau2026q11}.
    Given two such splittings with relative capacity load $\beta$, the unique proper
    orthochronous isometry of the reconstructed flat interval that maps one temporal axis
    to the other is the standard Lorentz boost
    \[
      \Lambda(\beta)
      =
      \begin{pmatrix}
        \gamma & -\gamma\beta \\
        -\gamma\beta & \gamma
      \end{pmatrix}.
    \]
    Length contraction follows as the reciprocal spatial projection onto the simultaneity
    hyperplane of the moving splitting:
    \[
      \ell = \frac{\ell_0}{\gamma}.
    \]
    The Born--Infeld capacity sphere parameterises admissible capacity splittings and is
    not a Lorentzian structure; the Lorentz group acts as the isometry group of the
    effective metric $\eta$, not of the BI sphere.
    Special-relativistic observer transformations are thereby reconstructed from the same
    capacity data that produced projective time dilation.
  \end{abstract}

  \medskip
  \noindent\textbf{Keywords.}
  Lorentz transformation; length contraction; inertial splitting; projective temporal
  ordering; Born--Infeld capacity sphere; effective Lorentzian metric; Cosmochrony.
  \enlargethispage{1\baselineskip}

  \newpage
  \tableofcontents

  \newpage
  \section{Introduction}
    \label{sec:introduction}

    The Lorentz-capacity sub-programme derives relativistic kinematics from the bounded
    Born--Infeld capacity budget of admissible projection on
    $\mathrm{Heis}_3(\mathbb{Z}/q\mathbb{Z})$.
    Paper~\cite{Beau2026n} identifies the projective temporal capacity $F^\tau = \sqrt{1-B^2}$
    as the Lorentz factor residual, and~\cite{Beau2026tp} proves that the temporal residual map
    \[
      P^{Q5b}_\tau \Sigma^{\mathrm{tot}}_n = \sqrt{1-B_n^2}
    \]
    is the unique positive function compatible with the Lorentzian splitting of Q5b--Q11,
    the Born--Infeld unit capacity sphere, and the Weil filling law.
    This closes time dilation ($d\tau/dn = 1/\gamma$) and the light cone ($B = 1
  \Rightarrow F^\tau = 0$).

    The open problem LC-O1 is stated in Section~7.2 of~\cite{Beau2026tp}:
    length contraction requires the reconstruction of effective Lorentz transformations
    between inertial observers from projective temporal ordering data.

    The resolution rests on a strict separation of two structures.
    The Born--Infeld capacity sphere
    \[
      (F^\tau)^2 + \beta^2 = 1
    \]
    is a Euclidean capacity relation.
    It parameterises admissible temporal splittings and determines $\beta$ and $1/\gamma$,
    but it is not the Minkowski interval.
    The Lorentzian metric is supplied by the homogeneous Q5b--Q11 closure:
    \[
      g^{\mu\nu} = 2\eta^{\mu\nu}.
    \]
    Lorentz transformations are the isometries of $\eta$, not symmetries of the BI sphere.
    The BI sphere determines which splittings are admissible; the Lorentz group determines
    how to pass between them.

    Section~\ref{sec:splittings} defines inertial observers as unit timelike splittings of
    the effective metric and identifies the rapidity with the projective capacity load.
    Section~\ref{sec:boost} proves that the unique interval-preserving isometry mapping
    one splitting to another is the Lorentz boost.
    Section~\ref{sec:contraction} derives length contraction.
    Section~\ref{sec:pto} states the connection to projective temporal ordering.
    Sections~\ref{sec:discussion} and~\ref{sec:conclusion} summarise the logical structure.

  \section{Inertial splittings and capacity load}
    \label{sec:splittings}

    Let $\Meff$ be the homogeneous effective spacetime reconstructed from the admissibility
    operator.
    The Q5b--Q11 closure gives the completed effective co-metric
    \[
      g^{\mu\nu} = 2\eta^{\mu\nu},
      \qquad
      \eta = \mathrm{diag}(-1,+1,+1,+1),
    \]
    with no free parameter~\cite{Beau2026q5b,Beau2026q11}.
    The corresponding metric is $g_{\mu\nu} = \tfrac{1}{2}\eta_{\mu\nu}$.
    The isometry group of $g$ is $\SO(3,1)$, the same as that of $\eta$.

    \begin{definition}[Inertial splitting]
      \label{def:inertial-splitting}
      An inertial observer is a future-directed unit timelike vector $u \in T_x\Meff$
      satisfying $\eta(u,u) = -1$, together with its orthogonal simultaneity hyperplane
      \[
        \Sigma_u := \bigl\{ X \in T_x\Meff \mid \eta(X,u) = 0 \bigr\}.
      \]
    \end{definition}

    \begin{remark}[Capacity sphere versus Lorentzian structure]
      \label{rem:bi-vs-lorentz}
      The Born--Infeld capacity sphere $(F^\tau)^2 + \beta^2 = 1$ is a Euclidean norm
      constraint on the capacity flow vector, not a Lorentzian structure.
      It provides the values of $\beta$ and $1/\gamma$ from projective data, but it does not
      determine the Lorentzian interval.
      The Lorentzian interval is supplied by $\eta$ from the Q5b--Q11 geometric
      closure~\cite{Beau2026q5b,Beau2026q11}.
      The Lorentz group acts as isometries of $\eta$; it does not act as isometries of the
      BI sphere.
    \end{remark}

    The rest observer corresponds to the splitting $u_0 = \partial_\tau$.
    An observer with normalised spatial capacity load $\beta$ in the $x$-direction has
    splitting
    \[
      u_\beta = \gamma(\partial_\tau + \beta\,\partial_x),
      \qquad
      \gamma := \frac{1}{\sqrt{1-\beta^2}}.
    \]
    One verifies directly that $u_\beta$ is unit timelike:
    $\eta(u_\beta, u_\beta) = -\gamma^2 + \gamma^2\beta^2 = -1$.

    The parameter $\beta$ is not introduced independently.
    By the Lorentz-capacity result of~\cite{Beau2026n,Beau2026tp}, the projective temporal
    capacity satisfies
    \[
      F^\tau = \sqrt{1-\beta^2} = \frac{1}{\gamma},
    \]
    so the rapidity of the inertial splitting is determined by the projective temporal
    capacity already computed in~\cite{Beau2026tp}.

  \section{Boost reconstruction}
    \label{sec:boost}

    Let $\xi$ be the rapidity defined by
    \[
      \beta = \tanh\xi, \qquad \gamma = \cosh\xi, \qquad \gamma\beta = \sinh\xi.
    \]
    Then $u_\beta = \cosh\xi\,\partial_\tau + \sinh\xi\,\partial_x$.

    \begin{lemma}[Uniqueness of the inertial boost]
      \label{lem:boost-uniqueness}
      There is a unique proper orthochronous linear isometry of $(T_x\Meff, \eta)$ that
      maps $u_0 = \partial_\tau$ to $u_\beta$ and acts as the identity on the transverse
      directions $\partial_y, \partial_z$.
      It is the Lorentz boost of rapidity $\xi$, acting on the $(\tau,x)$-plane by
      \[
        \begin{pmatrix} \partial_{\tau'} \\ \partial_{x'} \end{pmatrix}
        =
        \begin{pmatrix} \cosh\xi & \sinh\xi \\ \sinh\xi & \cosh\xi \end{pmatrix}
        \begin{pmatrix} \partial_\tau \\ \partial_x \end{pmatrix}.
      \]
    \end{lemma}

    \begin{proof}
      A proper orthochronous linear isometry of $\eta$ preserving $\partial_y$ and
      $\partial_z$ acts only on the $(\tau,x)$-plane.
      Preservation of the restriction $-d\tau^2 + dx^2$ requires a hyperbolic rotation:
      the isometry is of the form displayed above for some $\xi$.
      The condition $\partial_{\tau'} = u_\beta$ uniquely fixes
      $\cosh\xi = \gamma$ and $\sinh\xi = \gamma\beta$,
      hence $\xi = \mathrm{arctanh}\,\beta$.
    \end{proof}

    In coordinate form the boost reads
    \[
      \tau' = \gamma(\tau - \beta x), \qquad
      x'   = \gamma(x  - \beta\tau), \qquad
      y'   = y,                      \qquad
      z'   = z.
    \]

    \begin{theorem}[Lorentz transformations from projective temporal ordering]
      \label{thm:lorentz-from-pto}
      Given a projective temporal capacity load $\beta$ determined by~\cite{Beau2026n,Beau2026tp},
      the unique proper orthochronous transformation between the corresponding inertial
      splittings is the Lorentz boost
      \[
        \tau' = \gamma(\tau - \beta x),\qquad
        x'   = \gamma(x  - \beta\tau),\qquad
        y'   = y,\qquad z' = z,
      \]
      with $\gamma = (1-\beta^2)^{-1/2}$.
    \end{theorem}

    \begin{proof}
      The projective temporal capacity result of~\cite{Beau2026n,Beau2026tp} gives
      $F^\tau = (1-\beta^2)^{1/2} = 1/\gamma$, fixing $\beta$ and hence $\xi$.
      The Q5b--Q11 closure gives the effective co-metric $g^{\mu\nu} = 2\eta^{\mu\nu}$.
      Since the factor $2$ is a global conformal rescaling, the isometry group of $g$ is
      $\SO(3,1)$, the isometry group of $\eta$.
      By Lemma~\ref{lem:boost-uniqueness}, the unique proper orthochronous isometry of $\eta$
      mapping $\partial_\tau$ to $u_\beta$ and fixing the transverse directions is the
      displayed Lorentz boost.
    \end{proof}

  \section{Length contraction}
    \label{sec:contraction}

    \begin{corollary}[Length contraction]
      \label{cor:length-contraction}
      A rod of proper length $\ell_0$ at rest in the unprimed splitting has length
      $\ell = \ell_0/\gamma$ when measured on the simultaneity hyperplane of the primed
      splitting.
    \end{corollary}

    \begin{proof}
      Let the rod endpoints in the rest frame be $A = (\tau_A, x_A)$ and
      $B = (\tau_B, x_B)$ with $\Delta x = x_B - x_A = \ell_0$.
      The moving observer measures length at equal primed time, $\Delta\tau' = 0$.
      The boost gives $\Delta\tau' = \gamma(\Delta\tau - \beta\Delta x)$,
      so simultaneous in the primed frame means $\Delta\tau = \beta\Delta x$.
      The primed spatial separation is then
      \[
        \Delta x'
        = \gamma(\Delta x - \beta\Delta\tau)
        = \gamma(1-\beta^2)\Delta x
        = \frac{\Delta x}{\gamma}
        = \frac{\ell_0}{\gamma}.
      \]
    \end{proof}

  \section{Connection to projective temporal ordering}
    \label{sec:pto}

    The projective temporal ordering observable~\cite{Beau2026pto} supplies the scalar
    clock factor
    \[
      \frac{d\tau}{dn} = F^\tau = \frac{1}{\gamma}.
    \]
    LC-O1 shows that this factor is not an isolated time-dilation effect.
    Once the effective Lorentzian interval is reconstructed, the same capacity load $\beta$
    determines the rapidity $\xi = \mathrm{arctanh}\,\beta$ and hence the complete inertial
    transformation.

    Projective temporal ordering determines the clock axis of the observer.
    Q5b--Q11 determine the Lorentzian interval on which this axis lives.
    The boost is the unique interval-preserving change of admissible inertial splitting.

  \section{Discussion}
    \label{sec:discussion}

    The construction separates three layers, each with a distinct source.

    \textit{Layer 1: Born--Infeld capacity sphere.}
    The bounded BI admissibility budget of~\cite{Beau2026c} gives the unit capacity
    relation
    \[
      (F^\tau)^2 + \beta^2 = 1.
    \]
    This is a Euclidean norm constraint on the capacity flow vector.
    It is not the Minkowski interval.
    It provides $\beta$ and $1/\gamma$ from projective data.

    \textit{Layer 2: Effective Lorentzian co-metric.}
    The homogeneous geometric closure of Q5b--Q11~\cite{Beau2026q5b,Beau2026q11}
    gives the effective co-metric
    \[
      g^{\mu\nu} = 2\eta^{\mu\nu}.
    \]
    This supplies the interval whose isometry group is $\SO(3,1)$.
    The global factor $2$ is physically irrelevant for the isometry group.

    \textit{Layer 3: Inertial observer and boost.}
    An inertial observer is a unit timelike splitting of $\eta$.
    The rapidity of a splitting is fixed by the projective capacity load $\beta$ from
    Layer~1.
    The Lorentz boost is the unique isometry of $\eta$ (Layer~2) that maps one admissible
    splitting to another at capacity load $\beta$.

    Length contraction is a consequence of this boost, not an independent capacity law.
    It follows from the difference in simultaneity hyperplanes between two inertial
    splittings, given the interval structure of Layer~2.

    \begin{remark}[Logical status]
      \label{rem:logical-status}
      All inputs are previously closed: time dilation and $F^\tau = 1/\gamma$
      in~\cite{Beau2026n,Beau2026tp}; the effective co-metric $g^{\mu\nu} = 2\eta^{\mu\nu}$
      in~\cite{Beau2026q5b,Beau2026q11}; the BI admissibility budget
      in~\cite{Beau2026c,Beau2026m}.
      LC-O1 adds only the identification of inertial observers with unit timelike splittings
      and the application of the standard uniqueness of isometries of $\eta$ in each rapidity
      class.
      No new spectral or projective hypothesis is introduced.
    \end{remark}

  \section{Conclusion}
    \label{sec:conclusion}

    LC-O1 is closed.
    Projective temporal ordering data determine the capacity load $\beta$ and the residual
    $F^\tau = 1/\gamma$.
    The Q5b--Q11 closure supplies the flat effective Lorentzian interval with isometry
    group $\SO(3,1)$.
    An inertial observer is a unit timelike splitting of this interval; changing observer is
    changing the splitting while preserving the interval.
    The unique proper orthochronous isometry at capacity load $\beta$ is the Lorentz boost:
    \[
      \tau' = \gamma(\tau - \beta x), \qquad
      x'   = \gamma(x  - \beta\tau).
    \]
    Consequently, a rest length $\ell_0$ is measured as $\ell = \ell_0/\gamma$.

    The full Lorentz transformation, including length contraction, is thereby reconstructed
    from projective temporal ordering data and the effective metric normalisation.
    The Born--Infeld capacity sphere provides $\beta$ and $\gamma$;
    it is not the Minkowski interval and does not carry the Lorentz group as an isometry.

    \newpage
    \bibliographystyle{plain}
    \bibliography{cosmochrony-bibliography}

\end{document}
